%% ============================================================
%%  Section: Descriptive Statistics
%%  Depositor Discipline in Modern Banking (2026-04-07 draft)
%%
%%  Changes from 20260406 version:
%%    - Removed all [@@@ ...] placeholders and addressed each one.
%%    - §4.2 (Characteristics) rewritten: removed specific numbers throughout;
%%      retained only the most salient differences (size, deposit structure).
%%      CRE claim corrected: CRE share IS significantly higher for high-CECL
%%      banks in the new table (1.727pp, p<0.05), opposite of what the prior
%%      draft stated.
%%    - §4.2 shortened by roughly one-third. Performance paragraph collapsed
%%      to two sentences; deposit structure paragraph simplified.
%%    - §4.3 (Predictors) shortened to four sentences; removed coefficient
%%      recitation and moved the key takeaway to the opening sentence.
%%
%%  References two tables and two figures defined in tables_figures.tex:
%%    tab:descriptive_sumstats_cecl  (tab01_sumstats_20260403)
%%    tab:cecl_predictors_cross_section  (tab02_cecl_predictors_20260403)
%%    fig:cecl_equity_distribution   (fig01_cecl_dist_20260403)
%%    fig:cecl_adoption_quarter_distribution  (fig02_adoption_dist_20260403)
%%
%%  NOTE: The prior desc-stats file used \label{sec:data} which conflicts
%%  with the data section. This file uses \label{sec:descriptive}.
%% ============================================================

\section{Descriptive Statistics}
\label{sec:descriptive}

In this section, we describe the cross-sectional distribution of the treatment
variable, document the concentration of adoption timing, and compare pre-adoption
characteristics of high- and low-CECL banks.

\subsection{Distribution of the CECL Adjustment and Adoption Timing}
\label{sec:descriptive:distribution}

Figure~\ref{fig:cecl_equity_distribution} plots the cross-sectional distribution
of \texttt{CECL Adj}, the Day-One retained-earnings charge scaled by pre-adoption
book equity, for the community banks in the estimation sample.  The
distribution is right-skewed.  The full-sample mean is 0.28 percent of equity and
the median is essentially zero: roughly half of community banks adopted CECL with
loan portfolios whose forward-looking expected credit losses, under the new
standard, required little incremental reserve above what the incurred-loss
standard had already accumulated.  The right tail is substantial, however.  The
90th percentile of the distribution lies at approximately 2.5 percent of equity,
the threshold for our binary High-CECL indicator.  The banks that exceed
this threshold show a mean adjustment of nearly 3 percent of equity,
representing a material reduction in reported book capital from a single
mandatory accounting event.

Figure~\ref{fig:cecl_adoption_quarter_distribution} shows the distribution of
adoption quarters across the eligible six-quarter window (2022Q3--2023Q4).
Adoption is sharply concentrated in 2023Q1: 3,997 of the 4,320 sample banks
(92.5 percent) appear in their first CECL-compliant Call Report in that quarter,
exactly as prescribed for calendar-year non-public filers.  The remaining
323 banks are distributed across the five adjacent quarters, reflecting
non-calendar fiscal years, early voluntary adoption, and minor filing lags.
This near-uniform adoption calendar strengthens the identification: the
information shock reaches the market in a single concentrated quarter, making
it implausible that the timing itself was chosen in response to bank-specific
deposit conditions.  We nonetheless use bank-specific event time throughout,
assigning each bank its own $t^{*}$, so that the small minority with off-cycle
adoption dates contribute symmetric pre- and post-adoption windows to the panel.

\subsection{Characteristics of High- and Low-CECL Banks}
\label{sec:descriptive:balance}

Table~\ref{tab:descriptive_sumstats_cecl} compares High-CECL and Low-CECL banks
at the adoption cross-section.

High-CECL banks enter the adoption quarter with weaker loan quality and thinner
capital buffers.  Their nonperforming loan ratios and allowance ratios are both
significantly higher, and their CRE concentration is modestly but significantly
greater, consistent with more credit-sensitive loan portfolios requiring larger
expected-loss reserves under the new standard.  Capitalization is meaningfully
lower for high-CECL banks, reflecting the combined effect of elevated credit
exposure and the retained-earnings charge itself.

The pre-adoption deposit structure of the two groups is broadly similar on
dimensions central to our outcome variables.  Neither uninsured time deposits as
a share of assets nor the total uninsured deposit share differs significantly
between the groups.  One exception: High-CECL banks hold somewhat more insured
time deposits at baseline, reflecting greater reliance on retail interest-bearing
funding.  Bank fixed effects absorb this level difference, so our post-adoption
estimates capture within-bank changes.  The event-study pre-period coefficients
in Section~\ref{sec:results} confirm that deposit trends were parallel across
groups in the ten quarters before adoption (2020Q3 through 2022Q4).

Performance differences are modest in magnitude and absorbed by bank fixed
effects; they confirm that High-CECL banks represent a coherent group for which
the CECL disclosure carries the most new information about forward-looking credit
risk.

\subsection{What Predicts the CECL Adjustment?}
\label{sec:descriptive:predictors}

Table~\ref{tab:cecl_predictors_cross_section} regresses \texttt{CECL Adj} on
bank characteristics at the adoption quarter.  The adjustment is driven by
credit-risk fundamentals---NPL ratios and allowance levels---not by the deposit
composition or securities portfolio characteristics that might independently
generate depositor responses.  Unrealized securities losses, the channel central
to the 2023 stress at larger institutions, show no significant relationship with
the CECL adjustment in our community bank sample.  The overall fit is low
(adjusted $R^2$ of approximately one percent in the most saturated specification),
consistent with the CECL charge conveying information not easily recoverable from
pre-adoption observables.  Log assets is the strongest predictor across all
specifications, consistent with larger community banks holding more complex or
longer-duration loan portfolios that require greater lifetime-loss reserves
under the new standard.  Because bank fixed effects in the main regressions
absorb all time-invariant differences in size, any correlation between asset
scale and deposit dynamics does not confound the within-bank treatment
estimates.
